


\documentclass{class}

\usepackage{bbm}
\usetikzlibrary{positioning}


\begin{document}

\paragraph{Implementation by an ex ante nonlinear tariff and block-tariff approximation.}
Fix a category \(i\). In the direct mechanism of the paper, a type \(\theta\) consumer in state \(s\) receives quantity \(q_i(\theta,s)\), pays transfer \(t_i(\theta,s)\), and obtains utility
\[
\theta U\left(q_i(\theta,s),s\right)-t_i(\theta,s).
\]
Under incentive compatibility, the associated ex ante consumer surplus satisfies
\[
\E CS_i(\theta)
:=
\E_s\left[\theta U\left(q_i(\theta,s),s\right)-t_i(\theta,s)\right]
=
\E\underline{CS}_i
+
\int_{\underline{\theta}_i}^{\theta}
\E_s\left[U\left(q_i\left(\tilde{\theta},s\right),s\right)\right]\,d\tilde{\theta}.
\]
Equivalently, the expected transfer paid by type \(\theta\) is
\[
\E_s\left[t_i(\theta,s)\right]
=
\E_s\left[\theta U\left(q_i(\theta,s),s\right)\right]
-
\int_{\underline{\theta}_i}^{\theta}
\E_s\left[U\left(q_i\left(\tilde{\theta},s\right),s\right)\right]\,d\tilde{\theta}
-
\E\underline{CS}_i.
\]
Thus, \(\E\underline{CS}_i\) is the ex ante boundary term of the mechanism.

Because the optimal allocation \(q_i(\theta,s)\) is pointwise incentive compatible, it can be implemented by a state-contingent nonlinear tariff under which, after observing \(s\), the consumer chooses quantity directly:
\[
\max_{q\geq 0}\;\theta U(q,s)-\tau_i(q,s).
\]
Following the same logic as in Spulber, the tariff can be written directly with the boundary term as the fixed component:
\[
\tau_i(q,s)
=
-\E\underline{CS}_i
+
\int_0^q p_i(x,s)\,dx.
\]
The fixed term of the implementing tariff does not affect the consumer's choice once the state \(s\) is realized, since only the marginal price \(\partial_q\tau_i(q,s)\) enters the quantity-choice problem. 

Whenever \(\theta\mapsto q_i(\theta,s)\) is strictly increasing, let \(\theta_i(q,s)\) denote its inverse. The marginal price of the implementing nonlinear tariff is then
\[
p_i(q,s):=\partial_q\tau_i(q,s)=\theta_i(q,s)\,u(q,s).
\]
Hence the marginal price is equal to the willingness to pay of the type who selects quantity \(q\) in state \(s\).

Using the optimality condition from the planner's problem,
\[
u\left(q_i(\theta,s),s\right)\mathcal{G}_i(\theta)=\varepsilon_i(s),
\]
the marginal price charged at the allocation of type \(\theta\) can also be written as
\[
p_i\left(q_i(\theta,s),s\right)
=
\theta\,u\left(q_i(\theta,s),s\right)
=
\theta\,\frac{\varepsilon_i(s)}{\mathcal{G}_i(\theta)}.
\]
This expression makes the link with the primitives of the model transparent. Off peak, \(\varepsilon_i(s)=0\), so the marginal price is zero. On peak, the marginal price is positive and reflects both scarcity, through \(\varepsilon_i(s)\), and the incentive distortion, through \(\mathcal{G}_i(\theta)\).

A block tariff is obtained by discretizing this nonlinear tariff over a partition of the quantity space. Fix cutoffs
\[
0=q_{i,0}<q_{i,1}<\cdots<q_{i,B}.
\]
The economically correct incremental block price on block \(b\) is
\[
\pi_{i,b}(s)
:=
\frac{\tau_i\left(q_{i,b},s\right)-\tau_i\left(q_{i,b-1},s\right)}
{q_{i,b}-q_{i,b-1}}
=
\frac{1}{q_{i,b}-q_{i,b-1}}
\int_{q_{i,b-1}}^{q_{i,b}} p_i(x,s)\,dx.
\]
Because the fixed term cancels, each block price is simply the average marginal price over the block.

The associated incremental block tariff is
\[
\tau_i^{BT}(q,s)
=
-\E\underline{CS}_i
+
\sum_{b=1}^B
\pi_{i,b}(s)
\left(\min\left\{q,q_{i,b}\right\}-q_{i,b-1}\right)_+.
\]
By construction,
\[
\tau_i^{BT}\left(q_{i,b},s\right)=\tau_i\left(q_{i,b},s\right)
\qquad \text{for all } b=0,\dots,B.
\]
Thus, the block tariff preserves the payment increments of the nonlinear tariff while replacing its continuous marginal-price schedule by a step function. As the partition becomes finer, the block tariff converges to the original nonlinear tariff.

The shape of the block prices inherits the shape of the marginal-price schedule. If \(p_i(q,s)\) increases with \(q\), then the block prices \(\pi_{i,b}(s)\) increase across blocks; if \(p_i(q,s)\) decreases with \(q\), then the block prices decrease across blocks. In this sense, the block tariff is the natural discrete approximation of the nonlinear pricing rule that implements the optimal mechanism.


\paragraph{Endpoint behavior of the marginal price when $\beta_i = 0$.}
Fix a category $i$ and a peak state $s \in T_i$. When $\beta_i = 0$ the effective weight is $\mathcal{G}_i(\theta) = \Gamma_i(\theta)$, so the marginal price paid by type $\theta$ at the quantity she selects is $p_i(q_i(\theta,s),s) = \theta \cdot \varepsilon_i(s)/\Gamma_i(\theta)$.

At the lower bound, $J_i(\underline{\theta}_i) = \underline{\theta}_i - \gamma_i(\underline{\theta}_i)$ and $\Lambda_i(\underline{\theta}_i) = \gamma_i(\underline{\theta}_i)\mathbb{E}[\lambda_i(\tilde{\theta}) \mid \tilde{\theta} \geq \underline{\theta}_i] = \gamma_i(\underline{\theta}_i)\tilde{\lambda}_i$, so the two $\gamma_i$ terms cancel and $\Gamma_i(\underline{\theta}_i) = \tilde{\lambda}_i\underline{\theta}_i$. At the upper bound, $\gamma_i(\overline{\theta}_i) = 0$, so $J_i(\overline{\theta}_i) = \overline{\theta}_i$ and $\Lambda_i(\overline{\theta}_i) = 0$, giving $\Gamma_i(\overline{\theta}_i) = \tilde{\lambda}_i\overline{\theta}_i$.

In both cases $\theta/\Gamma_i(\theta) = 1/\tilde{\lambda}_i$, so the marginal price at both endpoints is the same:
\[
p_i\!\left(q_i(\underline{\theta}_i,s),s\right) = p_i\!\left(q_i(\overline{\theta}_i,s),s\right) = \frac{\varepsilon_i(s)}{\tilde{\lambda}_i}, \qquad s \in T_i.
\]
This is the implementation-space restatement of the boundary condition already established in the proof of Proposition~4: at both endpoints $\Gamma_i(\theta) = \tilde{\lambda}_i\theta$, which is why $R_{\mathcal{G}}(\underline{\theta}_i) = R_{\mathcal{G}}(\overline{\theta}_i)$ in the IR-unconstrained regime. Any nonlinearity of the tariff must therefore come from the interior behavior of $\theta/\Gamma_i(\theta)$, not from the endpoints.

\paragraph{Endpoint ranking of marginal prices when $\beta_i > 0$.}
Fix a category $i$ and a peak state $s \in T_i$. When $\beta_i > 0$ the effective weight is $\mathcal{G}_i(\theta) = \Gamma_i(\theta) + J_i(\theta)\beta_i$, so the marginal price paid by type $\theta$ at the quantity she selects is $p_i(q_i(\theta,s),s) = \theta \cdot \varepsilon_i(s)/\mathcal{G}_i(\theta)$.

At the upper bound, $\gamma_i(\overline{\theta}_i) = 0$ and $\Lambda_i(\overline{\theta}_i) = 0$, so $\Gamma_i(\overline{\theta}_i) = \tilde{\lambda}_i\overline{\theta}_i$ and $J_i(\overline{\theta}_i) = \overline{\theta}_i$, giving $\mathcal{G}_i(\overline{\theta}_i) = (\tilde{\lambda}_i + \beta_i)\overline{\theta}_i$ and therefore
\[
p_i\!\left(q_i(\overline{\theta}_i,s),s\right) = \frac{\varepsilon_i(s)}{\tilde{\lambda}_i + \beta_i}.
\]
At the lower bound, $J_i(\underline{\theta}_i) = \underline{\theta}_i - 1/g_i(\underline{\theta}_i)$ and $\Lambda_i(\underline{\theta}_i) = \gamma_i(\underline{\theta}_i)\tilde{\lambda}_i = \tilde{\lambda}_i/g_i(\underline{\theta}_i)$, so the two $\gamma_i$ terms cancel and $\Gamma_i(\underline{\theta}_i) = \tilde{\lambda}_i\underline{\theta}_i$ as before. Hence $\mathcal{G}_i(\underline{\theta}_i) = \tilde{\lambda}_i\underline{\theta}_i + \beta_i(\underline{\theta}_i - 1/g_i(\underline{\theta}_i)) = \underline{\theta}_i(\tilde{\lambda}_i + \beta_i - \beta_i/(\underline{\theta}_i g_i(\underline{\theta}_i)))$ and therefore
\[
p_i\!\left(q_i(\underline{\theta}_i,s),s\right) = \frac{\varepsilon_i(s)}{\tilde{\lambda}_i + \beta_i - \beta_i/(\underline{\theta}_i g_i(\underline{\theta}_i))}.
\]
Since $\beta_i/(\underline{\theta}_i g_i(\underline{\theta}_i)) > 0$, the denominator at the lower bound is strictly smaller than at the upper bound. Under Assumption~\ref{ass:interior}, $\mathcal{G}_i(\underline{\theta}_i) > 0$, so both prices are well defined and
\[
p_i\!\left(q_i(\underline{\theta}_i,s),s\right) > p_i\!\left(q_i(\overline{\theta}_i,s),s\right), \qquad s \in T_i.
\]
When $\beta_i > 0$, the lowest type therefore faces a strictly higher endpoint marginal price than the highest type. Unlike the case $\beta_i = 0$, where endpoint prices are equalized at $\varepsilon_i(s)/\tilde{\lambda}_i$, the binding budget constraint breaks this equality by reducing the effective weight at the bottom of the distribution through the $J_i(\underline{\theta}_i)\beta_i$ term, which is negative whenever $J_i(\underline{\theta}_i) < 0$.

\paragraph{Derivative of the implemented marginal price.}
Fix a category $i$ and a peak state $s \in T_i$. Write $q(\theta) := q_i(\theta,s)$, $\varepsilon := \varepsilon_i(s)$, and $\mathcal{G}(\theta) := \mathcal{G}_i(\theta)$. On peak, the optimal allocation satisfies $u(q(\theta),s)\mathcal{G}(\theta) = \varepsilon$, and the implemented marginal price is $p_i(q(\theta),s) = \theta\,\varepsilon/\mathcal{G}(\theta)$. Differentiating the first-order condition with respect to $\theta$ gives $q_\theta(\theta) = -\varepsilon\,\mathcal{G}'(\theta)/(u_q(q(\theta),s)\mathcal{G}(\theta)^2)$, and differentiating the price gives $\partial_\theta p_i(q(\theta),s) = \varepsilon(\mathcal{G}(\theta) - \theta\mathcal{G}'(\theta))/\mathcal{G}(\theta)^2$. Hence
\[
p_{i,q}\!\left(q_i(\theta,s),s\right) = \frac{\partial_\theta p_i(q(\theta),s)}{q_\theta(\theta)} = -\,u_q\!\left(q_i(\theta,s),s\right)\frac{\mathcal{G}_i(\theta) - \theta\mathcal{G}_i'(\theta)}{\mathcal{G}_i'(\theta)}.
\]

\section{The two regimes.}
When $\beta_i = 0$, $\mathcal{G}_i(\theta) = \Gamma_i(\theta) = \tilde{\lambda}_i J_i(\theta) + \Lambda_i(\theta)$, so the formula reduces to
\[
p_{i,q}\!\left(q_i(\theta,s),s\right) = -\,u_q\!\left(q_i(\theta,s),s\right)\frac{\Gamma_i(\theta) - \theta\Gamma_i'(\theta)}{\Gamma_i'(\theta)}.
\]
When $\beta_i > 0$, $\mathcal{G}_i(\theta) = \Gamma_i(\theta) + J_i(\theta)\beta_i$ and $\mathcal{G}_i'(\theta) = \Gamma_i'(\theta) + J_i'(\theta)\beta_i$. Substituting and expanding the numerator,
\[
\mathcal{G}_i(\theta) - \theta\mathcal{G}_i'(\theta) = \Gamma_i(\theta) - \theta\Gamma_i'(\theta) + \left(\theta\gamma_i'(\theta) - \gamma_i(\theta)\right)\beta_i,
\]
where we used $J_i(\theta) - \theta J_i'(\theta) = \theta\gamma_i'(\theta) - \gamma_i(\theta)$. Using $\Gamma_i(\theta) = \tilde{\lambda}_i J_i(\theta) + \Lambda_i(\theta)$ the formula becomes
\[
p_{i,q}\!\left(q_i(\theta,s),s\right) = -\,u_q\!\left(q_i(\theta,s),s\right)\frac{\left(\tilde{\lambda}_i + \beta_i\right)\left(\theta\gamma_i'(\theta) - \gamma_i(\theta)\right) + \Lambda_i(\theta) - \theta\Lambda_i'(\theta)}{\Gamma_i'(\theta) + J_i'(\theta)\beta_i}.
\]
Setting $\beta_i = 0$ recovers the previous case.

\end{document}

